\documentclass[letterpaper,12pt,oneside,openany, spanish, oldfontcommands]{memoir}

\usepackage[spanish,activeacute,es-noshorthands]{babel}
\usepackage[latin1]{inputenc}
\usepackage{setspace}
\usepackage{tabularx}
\usepackage{anysize}
\usepackage{algorithm}
\usepackage{algorithmic}
%archivo para que los algoritmos me salgan en español
\floatname{algorithm}{C\'odigo}
\renewcommand{\listalgorithmname}{INDICE DE C\'ODIGOS}
\renewcommand{\algorithmicrequire}{\textbf{Entrada:}}
\renewcommand{\algorithmicensure}{\textbf{Salida:}}
\renewcommand{\algorithmicend}{\textbf{fin}}
\renewcommand{\algorithmicif}{\textbf{si}}
\renewcommand{\algorithmicthen}{\textbf{entonces}}
\renewcommand{\algorithmicelse}{\textbf{si no}}
\renewcommand{\algorithmicelsif}{\algorithmicelse,\ \algorithmicif}
\renewcommand{\algorithmicendif}{\algorithmicend\ \algorithmicif}
\renewcommand{\algorithmicfor}{\textbf{para}}
\renewcommand{\algorithmicforall}{\textbf{para todo}}
\renewcommand{\algorithmicdo}{\textbf{hacer}}
\renewcommand{\algorithmicendfor}{\algorithmicend\ \algorithmicfor}
\renewcommand{\algorithmicwhile}{\textbf{mientras}}
\renewcommand{\algorithmicendwhile}{\algorithmicend\ \algorithmicwhile}
\renewcommand{\algorithmicloop}{\textbf{repetir}}
\renewcommand{\algorithmicendloop}{\algorithmicend\ \algorithmicloop}
\renewcommand{\algorithmicrepeat}{\textbf{repetir}}
\renewcommand{\algorithmicuntil}{\textbf{hasta que}}
\renewcommand{\algorithmicprint}{\textbf{imprimir}} 
\renewcommand{\algorithmicreturn}{\textbf{devolver}} 
\renewcommand{\algorithmictrue}{\textbf{cierto }} 
\renewcommand{\algorithmicfalse}{\textbf{falso }} 

\usepackage{array}
 %Inclusi\'on de gr\'aficos al lado de texto
\usepackage{wrapfig}
\usepackage{colortbl}
\usepackage{graphicx}
\usepackage[left=4cm,top=3cm,right=3cm,bottom=3cm]{geometry}
\usepackage{slashbox}
\usepackage{supertabular}
\usepackage{longtable}
\usepackage{rotating}
%para usar multiples filas
\usepackage{multirow}
%para usar multiples columnas
\usepackage{dcolumn}
\usepackage{pdflscape}
\usepackage{amsmath}
%evitar el reset de los footnotes entre cada cap'itulo
\usepackage{remreset}
\makeatletter\@removefromreset{footnote}{chapter}\makeatother
%footnotes en la parte de abajo de la pagina
\usepackage[bottom]{footmisc}
\usepackage{titlesec}
\usepackage{float}
%Incluir bibliograf\'ia en el \'indice
\usepackage{tocbibind} 
\usepackage{pdflscape}
\usepackage{amsmath}
% paquete para el uso de simbolos como el checkmark
\usepackage{bbding}
%Inclusi\'on de gr\'aficos al lado de texto
\usepackage{wrapfig} 
\usepackage{colortbl} 
\usepackage{titlesec}
\usepackage{float}
\usepackage{tabularx}
\usepackage{anysize}
\usepackage{array}
\usepackage{color}
%para los simbolos usados en las vinetas
\usepackage{pifont}
\usepackage[dvips,final]{epsfig}
\definecolor{vivinicered}{RGB}{0,104,166}
\definecolor{cap}{RGB}{204,230,255}
% Para tener marcadores por secciones y subsecciones
\usepackage[pdftex=true,colorlinks=true,plainpages=false]{hyperref}
\hypersetup{
    colorlinks,
    citecolor=vivinicered,
    filecolor=vivinicered,
    linkcolor=vivinicered,
    urlcolor=vivinicered,
}

%\usepackage[table,svgnames,hyperref]{xcolor}
\usepackage{calc}
\usepackage{ifthen}
\usepackage{pgf,tikz}
\usetikzlibrary{shapes,arrows,calc,positioning,shadows}

%%%%%%%%	COMO SE VEN LOS CAPITULOS	%%%%%%%%%%
\usepackage{color,calc,graphicx,soul,fourier}
\makeatletter
\newlength\dlf@normtxtw
\setlength\dlf@normtxtw{\textwidth}
\def\myhelvetfont{\def\sfdefault{mdput}}
\newsavebox{\feline@chapter}
\newcommand\feline@chapter@marker[1][4cm]{%
\sbox\feline@chapter{%
\resizebox{!}{#1}{\fboxsep=1pt%
\colorbox{vivinicered}{\color{white}\bfseries\sffamily\thechapter}%
}}%
\rotatebox{90}{%
\resizebox{%
\heightof{\usebox{\feline@chapter}}+\depthof{\usebox{\feline@chapter}}}%
{!}{\scshape\so\@chapapp}}\quad%
\raisebox{\depthof{\usebox{\feline@chapter}}}{\usebox{\feline@chapter}}%
}
\newcommand\feline@chm[1][4cm]{%
\sbox\feline@chapter{\feline@chapter@marker[#1]}%
\makebox[0pt][l]{% aka \rlap
\makebox[1cm][r]{\usebox\feline@chapter}%
}}
\makechapterstyle{daleif1}{
\renewcommand\chapnamefont{\normalfont\Large\scshape\raggedleft\so}
\renewcommand\chaptitlefont{\normalfont\huge\bfseries\scshape\color{vivinicered}}
\renewcommand\chapternamenum{}
\renewcommand\printchaptername{}
\renewcommand\printchapternum{\null\hfill\feline@chm[2.5cm]\par}
\renewcommand\afterchapternum{\par\vskip\midchapskip}
\renewcommand\printchaptertitle[1]{\chaptitlefont\raggedleft ##1\par}
}
\makeatother
\chapterstyle{daleif1}
%%%%%%%%%%%%%%%%%%%%%%%%%%%%%%%%%%%%%%%%%%%%%%%%%%%%%%%%%%%%%%%%%%%%%%%%%%%%%%%%%%%%%%%%%%%%%%%%

%evitar el reset de los footnotes entre cada cap'itulo
\usepackage{remreset}
\makeatletter\@removefromreset{footnote}{chapter}\makeatother

%footnotes en la parte de abajo de la pagina
\usepackage[bottom]{footmisc}

\setcounter{secnumdepth}{3}
\setcounter{tocdepth}{4}

\setcounter{secnumdepth}{3}
% el nivel de detalle del indice. En 3 salen hasta el subsubsection
\setcounter{tocdepth}{3}

% espaciado de 1.5 entre lineas
\renewcommand{\baselinestretch}{1.5}

\AtBeginDocument{
  \def\labelitemi{\ding{224}}
  \def\labelitemii{\ding{223}}
}

%%%%%%%%	COMANDOS PARA IMAGENES PNG	%%%%%%%%
% \pngfig
% Usado para introducir imagenes adentro del directorio img/ cuya extension sea .png
% Float -> valor es el porcentaje de tama~no, 
% String -> nombre del archivo (funciona con archivos con nombres sin espacios)
% String -> Descripci'on que va a ir debajo de la imagen introducida
% String -> Fuente de la figura
% \pngfig[.5]{archihbhjv{descripcion de la imagen introducida} 
\newcommand{\pngfig}[4][1.0]{
\begin{figure}[H]
\begin{centering}
\scalebox{#1}{\includegraphics{/Users/vivianatrujillo/tesis/plantilla/imagenes/#2.png}}
\par\end{centering}
\begin{centering} 
\caption{\footnotesize #3} 
\label{fig:#2}
\par\end{centering}
\centering{}{\tiny #4}
\end{figure}
}

\newcommand{\pngfigrotada}[4][1.0]{
 \begin{figure}[H]
 \centering
 \vspace{2mm}
 \scalebox{#1}{\begin{sideways} \includegraphics{/Users/vivianatrujillo/tesis/plantilla/imagenes/#2.png} \end{sideways} }
 \caption{\footnotesize #3}
 \label{fig:#2}
{\tiny
 #4}
 \end{figure}
}
%%%%%%%%	COMANDOS PARA IMAGENES JPG	%%%%%%%%
\newcommand{\jpgfig}[4][1.0]{
\begin{figure}[H]
\begin{centering}
\scalebox{#1}{\includegraphics{/Users/vivianatrujillo/tesis/plantilla/imagenes/#2.jpg}}
\par\end{centering}
\begin{centering}
\begin{quote} 
\caption{\footnotesize #3}
\label{fig:#2}
\part\end{centering}
\centering{}{\tiny #4}
\end{quote}
\end{figure}
}

% Igual que las anteriores pero esta vez coloca la imagen rotada 90 grados (landscape)
\newcommand{\jpgfigrotada}[4][1.0]{
 \begin{figure}[H]
 \centering
 \vspace{2mm}
 \scalebox{#1}{\begin{sideways} \includegraphics{/Users/vivianatrujillo/tesis/plantilla/imagenes/#2.jpg} \end{sideways} }
 \caption{\footnotesize #3}
 \label{fig:#2}
{\tiny
 #4}
 \end{figure}
}

%%%%%%%%	COMANDOS PARA TABLAS	%%%%%%%%
% \nuevatabla
% Inserta una nueva tabla centrada y con los atributos.
% String -> Nombre de la tabla para luego ser referenciada (Sin espacios)
% String -> Texto descriptivo de la tabla que aparecer'a debajo de la misma.
% String -> Fuente elaboraci'on de la tabla
% Opciones de las columnas introducidas
% EJEMPLO: \nuevatabla{resumenDecripcionesAtributosCalidad}{Resumen de descripciones de atributos de calidad}{Fuente: Elaboraci'on propia}{|c|p{10cm}|}{\textbf{Atributo de Calidad} & \textbf{Descripci'on} \\ \hline}
% Ojo, recordar que las opciones de posici'on de la tabla se est'an colocando por defecto desde este comando y no son escogidos en la invocaci'on. Recordar qu'e:
%b: Intentara ponerla en el fondo de la pagina
%h: Intentara ponerla en la misma posicion en que se encuentra en el codigo fuente
%t: Intentara ponerla al principio de la pagina
%p: La pondra en una pagina que contenga solo elementos flotantes
%!: Ignorar la mayora de las restricciones impuestas por LATEX
\newcommand{\nuevatabla}[5]{
{\scriptsize 
\begin{table}[H] 
\setlength{\arrayrulewidth}{1pt}
\begin{center}
\begin{tabular}{#4}\hline
#5
\hline
\end{tabular}
\begin{quote} 
\caption{\footnotesize #2}
\label{tbl:#1}
\begin{center}
{\tiny #3}
\end{center}
\end{quote}
\end{center}
\end{table}
}
}

% lo mismo que la anterior pero con la tabla rotada 90 grados
\newcommand{\nuevatablarotada}[5]{
\begin{landscape}
\begin{table}[H] \footnotesize %[!hbt]%[htdp]
\setlength{\arrayrulewidth}{1pt}
\begin{center}
\caption{\footnotesize #2}
{\tiny
#3\\
\begin{tabular}{#4}\hline
#5}
\label{tbl:#1}
\end{tabular}
\end{center}
\end{table}
\end{landscape}
}

% \nuevalongtabla
%Cuando la tabla es muy grande y no cabe en una p'agina es necesario el uso de esta funci'on
% es igual a la \nuevatabla, pero garantiza que se pueda pasar a m'as de una p'agina.
%1 label
%2 caption de la tabla
%3 fuente
%4 la estructura
%5 t'itulo que se debe mantener en otras hojas
%6 contenido
\newcommand{\nuevalongtabla}[6]{
\begin{longtable}[H] 	 {#4}
\hline
\endfirsthead
\hline
#5 \\ 
\hline 
\endhead 
#6
\hline 
\caption{\footnotesize #2 - #3}
\label{tbl:#1}
\end{longtable}
}

% \nuevalongtablarotada
%Cuando la tabla es muy grande y no cabe en una p'agina es necesario el uso de esta funci'on
% es igual a la \nuevatabla, pero garantiza que se pueda pasar a m'as de una p'agina.
%1 label
%2 caption de la tabla
%3 fuente
%4 la estructura
%5 t'itulo que se debe mantener en otras hojas
%6 contenido
\newcommand{\nuevalongtablarotada}[6]{
\begin{landscape}
\begin{longtable}[H]{#4} 
\caption{\footnotesize #2}
#3\\
\hline
\endfirsthead
\hline
#5 \\ 
\hline 
\endhead 
#6
\label{tbl:#1}
\end{longtable}
\end{landscape}
}

%%%%%%%%%%%%%%%%%%%%%%%%%%%%%%%%%%%%%%%%%%%% COMIENZO DEL DOCUMENTO %%%%%%%%%%%%%%%%%%%%%%%%%%%%%%%%%%%%%%%%%%%%%%%%%%%%%%%%%%%%%%

\begin{document}

\renewcommand{\contentsname}{'INDICE DE CONTENIDO}
\renewcommand{\bibname}{Bibliograf\'ia} % Para que ya no diga Bibliograf'ia sino diga Referencias
\renewcommand{\tablename}{Tabla} % Los Captions de tablas salen con la palabra Tabla y no con la palabra Cuadro
\renewcommand{\listtablename}{'INDICE DE TABLAS} % Para que diga 'indice de tablas y no 'indice de cuadros
\renewcommand{\listfigurename}{'INDICE DE FIGURAS}
\renewcommand{\chaptername}{CAP'ITULO}
    
\parskip 6pt

\pagenumbering{roman}
\pagestyle{empty}

\begin{center}
\parbox{4cm}{\includegraphics[width=0.16\textwidth]{/home/israelord/Dissertation/Book/my_files/images/escudoUcab.jpg}}\parbox{9cm}{\begin{center}
UNIVERSIDAD CATÓLICA ANDRÉS BELLO \\
FACULTAD DE INGENIERÍA \\
ESCUELA DE INGENIERÍA INFORMÁTICA \\
\end{center}}

\vspace{\fill}


\begin{LARGE}
Desarrollo de una ontología y un buscador utilizando tecnologías de Web Semántica
\end{LARGE}

\vspace{\fill}

TRABAJO ESPECIAL DE GRADO\\
presentado ante la\\
UNIVERSIDAD CATOLICA ANDRES BELLO\\
como parte de los requisitos para optar al t\'itulo de\\
INGENIERO EN INFORM\'ATICA

\vspace{\fill}

{
\begin{tabular}{cl}
REALIZADO POR & \textbf{Israel Fermín Montilla}\tabularnewline
\multicolumn{2}{c}{}\tabularnewline
PROFESOR GUIA & \textbf{Wilmer Pereira}\tabularnewline
\multicolumn{2}{c}{}\tabularnewline
FECHA & Caracas, septiembre de 2011\tabularnewline
\end{tabular}
}


\end{center}

\chapter*{AGRADECIMIENTOS}
\addcontentsline{toc}{chapter}{AGRADECIMIENTOS} 

...

Me gustaría agradecer, en primer lugar, a ese ser superior al que algunos llaman Dios, otros Alá, otros Buda, otros Krishna y que muchos, simplemente, no sabemos cómo llamarlo. A mi familia, que siempre me ha apoyado en mis metas y mis locuras. A mi papá y a mi mamá por estar siempre pendientes de mi formación académica y personal. A Giselle por todo su cariño, su apoyo y por siempre creer en mi y por hacer que mi último semestre fuera inolvidable. A mi tutor, el Prof. Wilmer Pereira, y a mi amigo el Ing. Carlos Pérez, por su apoyo incondicional durante la realización de este trabajo a nivel práctico y a nivel conceptual. Sin ustedes, nada de esto habría sido posible.

También me gustaría agradecer a quienes me acompañaron semestre a semestre, mis amigos Gerardo Barcia, Jonathan Trujillo, Khaterine Castellano, Viviana Trujillo, Ronald Oribio y Freddy Florenca. A los profesores Carlos Barroeta, Lúcia Cardoso, Rodolfo Campos y Darío León, sin sus enseñanzas me sentiría inseguro trabajando en la calle en cualquier área de la carrera. Al Prof. Ricardo Casanova por brindarme su amistad durante los últimos semestres y a mis alumnos de preparaduría, especialmente a Marcos Ackerman, Ayleen Posadas, Héctor Sam y Juan Francisco Perozo, siento que fue más lo que ustedes me dejaron a mi que lo que yo dejé en ustedes.

A todos quienes, de alguna u otra manera me ayudaron,
Muchas Gracias.

...

\begin{flushright}
\textit{Israel Fermín Montilla}
\end{flushright}


\pagestyle{plain}

\begin{SingleSpace}
\tableofcontents

\newpage

\listoffigures

\newpage

%\listoftables

%\newpage

%\listofalgorithms
\end{SingleSpace}	


\pagenumbering{arabic}
\pagestyle{empty}

\makepagestyle{ruled}
\makeoddfoot{ruled}{}{}{\thepage}
\makeheadrule{ruled}{\textwidth}{\normalrulethickness}
\makeevenhead{ruled}{\scshape\leftmark}{}{} % small caps
\makeoddhead{ruled}{}{}{\rightmark}
\pagestyle{ruled}


\chapter*{SINOPSIS}
\addcontentsline{toc}{chapter}{RESUMEN} 

En el presente Trabajo Especial de Grado se ....
\chapter*{Introducci\'on}
\label{chap:introduccion}
\addcontentsline{toc}{chapter}{Introducción} 

Copiar y pegar lo mismo de la propuesta

\end{itemize}

\chapter{Problema}
\label{chap:problema}

%%%%%%%%%%PLANTEAMIENTO DEL PROBLEMA%%%%%%%%%%%%%%

\section{Planteamiento del Problema }

....

Copy-pasta de la propuesta

...

%%%%%%%%%%%OBJETIVOS%%%%%%%%%

\section{Objetivo General }

Copy-pastas de la propuesta

\section{Objetivos Espec\'ificos }

\begin{enumerate}
	\item ...
	\item ...
	\item ...
\end{enumerate}

%%%%%%%%%%%ALCANCE%%%%%%%%%%%

\section{Alcance }

Copy-Pasta de la propuesta

%%%%%%%%%LIMITACIONES%%%%%%%

\section{Limitaciones }
...
%%%%%%%%%JUSTIFICACION%%%%%%%

\section{Justificaci\'on }

...

\chapter{Marco Teórico}
\label{chap:marcoTeorico}

Para poder adentrarnos en el tema de la Web Semántica, resulta necesario estudiar y comprender los antecedentes y los eventos que poco a poco han ido llevando a la evolución de la Internet a la herramienta que tenemos hoy día y que, poco a poco también, irán llevando la Web actual, la Web 2.0, a su evolución natural: la Web 3.0, es decir, la Web Semántica.

Si actualmente nos encontramos en la llamada Web 2.0, tuvo que, en algún momento, existir una Web 1.0. Esta primera versión, por así llamarla, sólo contaba con portales que únicamente exponían contenido. La Web 1.0 estaba destinada a la publicación de contenidos corporativos, no daba la posibilidad de participación abierta a los usuarios, no existían espacios para la publicación de contenido por parte de los usuarios y, estos usuarios, eran importantes en tanto fueran consumidores, estas, según Ricardo Casanova (2010), eran algunas de las características principales de la Web 1.0 que, además, resultan ser grandes limitaciones. En la Web 1.0, únicamente personas especializadas eran capaces de crear contenido, por ello, sólo las grandes empresas podían disponer de un espacio en la red, el resto de los usuarios únicamente podían recibir (consumir) el mensaje o el contenido que era publicado, sin la posibilidad de participar en la generación y actualización del mismo.

El término Web 2.0 aparece a mediados del año 2004, y fue creciendo paulatinamente hasta convertirse en portada de los principales seminarios y congresos en la navidad de 2006. Según O'Reilly (citado por Pardo y Cobo, 2007), principal promotor de la Web 2.0, algunos de sus principios básicos son: La web como plataforma, el aprovechamiento de la Inteligencia Colectiva, la gestión de Bases de Datos como competencia básica, el software no limitado a un solo dispositivo y brindar experiencias enriquecedoras al usuario. Todo esto quiere decir, que ahora la interacción de los usuarios es bidireccional, sigue siendo un subgrupo técnico de estos usuarios el que crea los portales, pero ahora, la gran diferencia es que todos los usuarios son capaces de generar contenido en dichos portales. El usuario ya no es un simple consumidor, sino que pasa a ser consumidor-generador de contenido en la web. Pasamos de una red pasiva, de páginas estáticas, a una red activa de páginas dinámicas que interactúan con bases de datos para almacenar y actualizar su contenido que, a su vez, es generado por los usuarios que hacen vida dentro del portal. 

La evolución es un factor constante en todas las áreas y tecnologías de las Ciencias de la Computación y la Ingeniería Informática, Internet no es la excepción a esta regla, por lo es fácil deducir que la Web 2.0 no es más que el estado actual y transitorio de esta tecnología que, eventualmente, evolucionará a una Web 3.0, cuyas características están siendo definidas por el W3C bajo el marco de la Web Semántica.

\section{La Web Semántica}

El W3C (S/F) define a la Web Semántica como ``una Web extendida, dotada de mayor significado en la que cualquier usuario de Internet, podrá encontrar respuestas a sus preguntas de forma más rápida y sencilla gracias a que la información está mejor definida''. Al dotar a la Web de un mayor significado y, por consiguiente, de mayor semántica, es posible obtener solución a problemas comunes de búsqueda de información, gracias a la implementación de una infraestructura y lenguajes comunes de búsqueda, destinados a resolver dichos problemas, realizando dichas búsquedas tomando en cuenta el contexto y el significado real de la consulta.

Una vez definida la Web Semántica a nivel conceptual, conviene examinar brevemente cómo funciona. Supongamos que la Web tiene la capacidad de construir una base de conocimiento sobre las preferencias de los usuarios y que, a través de una combinación entre su conocimiento y la información disponible en la Internet, sea capaz de atender de forma exacta las demandas de información por parte de los usuarios, por ejemplo, reserva de hoteles, vuelos, médicos o libros.

Si esto ocurriese en la vida real, el usuario, al intentar encontrar ``todos los vuelos a Praga para mañana por la mañana'', obtendría resultados exactos sobre su búsqueda. Desafortunadamente la realidad es otra, los buscadores actuales mostrarían resultados acerca de ``Praga'' como localidad turística, noticias de sucesos ocurridos en ``Praga'', quizás páginas de periódicos locales, foros sobre ``Praga'', en resumen, sitios que contienen las palabras que conforman la consulta hecha por el usuario. Estos resultados son inexactos y, por sí solos, no satisfacen las necesidades de información del usuario, es necesario un segundo filtro, examinar uno a uno los resultados y extraer manualmente la información que resulte interesante. Por otro lado, un buscador con capacidad semántica, mostraría información más exacta a lo que se desea obtener. La ubicación geográfica desde la cual se envía la consulta sería detectada de manera automática, sin necesidad de indicar el punto de partida, además, elementos de la oración como ``mañana'' adquirirían significado, siendo convertidos en un día concreto calculado en función de un ``hoy''. De igual manera ocurriría con el segundo ``mañana'', el cual sería interpretado como un momento determinado del día. Todo ello a través de una Web en la que los datos, dejan de ser sólo datos y pasan a ser información llena de significado.

La manera como se procesará la información no sólo será en términos de entrada y salida de parámetros de búsqueda, sino en términos de su semántica, apoyándose en una infraestructura basada en meta-datos. Vale acotar, que no se trata de Inteligencia Artificial, sino de dar a las máquinas la capacidad de resolver problemas bien definidos, a través de operaciones bien definidas, ejecutadas sobre datos existentes bien definidos a través de meta-datos.

Dentro de la Web Semántica convergen una serie de estándares y tecnologías, muchas de ellas aún en proceso de definición y desarrollo por parte de investigadores alrededor del mundo y el W3C, las cuales se explican brevemente a continuación:

\section{Ontologías}

Según el W3C (S/F), ``una Ontología define los términos utilizados para describir y representar conocimiento en un área''. Las Ontologías son utilizadas por personas y aplicaciones que necesitan compartir conocimiento acerca de un área específica, es por ello que las ontologías son relativas a un tópico o área en especial, además, incluyen información acerca de los conceptos básicos y las relaciones entre ellos, esta información puede ser utilizada por las computadoras. De esta manera, una Ontología codifica el conocimiento de un área y lo hace accesible y reutilizable.

La palabra Ontología, se ha utilizado a lo largo de la historia para definir elementos con distintos grados de estructura: desde taxonomías simples hasta teorías lógicas complejas. La Web Semántica necesita ontologías con cierto grado de estructura y significado para poder especificar descripciones para los siguientes conceptos:

\begin{itemize}
\item Clases y dominios de interés.
\item Las relaciones que pueden existir entre las clases descritas.
\item Las propiedades o atributos que pueden tener las clases descritas.
\end{itemize}

Las Ontologías son escritas en lenguajes basados en lógica, de esta manera se tiene la garantía de que son precisas, detalladas, consistentes y significativas (W3C, S/F). Muchas herramientas para Ontologías pueden realizar procesos de razonamiento sobre el conocimiento que definen, de esta manera, se puede tener cierta ``inteligencia'' y se pueden desarrollar aplicaciones con capacidades complejas como consulta de información de manera semántica y conceptual, soporte a la toma de decisiones, gestión de conocimiento, bases de datos inteligentes, comercio electrónico y entendimiento del lenguaje natural.

En la Web Semántica, resulta necesario el uso de Ontologías ya que proporciona una manera sencilla y eficiente de representar la semántica de los documentos descritos y permitir que sean utilizadas y consultadas por aplicaciones y agentes de software. Para el W3C, el uso de ontologías, permitirá a las aplicaciones del futuro actuar de manera ``inteligente'' para cumplir con su trabajo de una manera más rápida y con mayor exactitud.

\section{RDF}

Sus siglas significan ``Resource Description Framework'', y es ``un lenguaje para la descripción de recursos disponibles en la World Wide Web'' (W3C, 2004). Es un modelo estándar para el intercambio de datos en la Web, RDF extiende la estructura de enlaces de la Web utilizando URIs para dar nombre a las cosas, es así como se establecen enlaces entre dos extremos: el sujeto y el objeto, a través de una relación o propiedad definida a través de RDF. Normalmente, a esta manera de expresar relaciones, se le conoce como ``triples''. Este modelo de relaciones permite que los datos estructurados o semi-estructurados sean mezclados, expuestos y compartidos a través de diversas aplicaciones en distintas plataformas.

Toda la estructura descrita en un RDF forma una estructura de grafo, donde los arcos representan las relaciones entre dos recursos, los cuales son representados por los nodos de dicho grafo (W3C, S/F). Este modelo mental, es la manera más utilizada para lograr explicaciones visuales y fáciles de comprender.

La sintaxis de RDF, está basada en XML, al igual que muchos otros metalenguajes y lenguajes de marcado, como HTML y XHTML. De esta manera, es posible además definir atributos adicionales a los recursos que son descritos en el documento.

\section{RDFS}

Siendo RDF un lenguaje basado en notación XML, pareciera lógico que siga más o menos el mismo esquema. Si recordamos, y si somos estrictos, cada XML debería tener un XML Schema (o, en su defecto, un DTD, pero estos están siendo poco a poco desplazados por los XMLS), que resulta ser otro archivo XML, que lo define. RFDS significa ``RDF Schema'' y es ``un lenguaje para la descripción de vocabularios en la Web'' (GSI, S/F). Es una extensión semántica de RDF para describir vocabulario, RDFS no busca definirlos, sino describirlos, proporcionar la facilidades para describir tipos y clases de un mismo dominio y servir como sistema de tipos para RDF.

Un RDF Schema, se escribe utilizando las mismas reglas de RDF, es decir, el Esquema RDF es un documento RDF que, a su vez, es capaz de definir a otros RDF. El RDFS, extiende al RDF tradicional para poder implementar jerarquías de clases y relaciones un poco más complejas que las que RDF permite definir por sí solo.

RDFS, al igual que RDF, aún se encuentran en proceso de definición y desarrollo, aunque ya el trabajo está bien adelantado y se puede trabajar y desarrollar con ambos. Un ejemplo de esto es el proyecto FOAF o ``Friend of a Friend'', que es un proyecto de la Web Semántica para describir personas y las relaciones entre ellas utilizando documentos RDF.

\section{OWL}

El W3C (2005), define OWL como ``un Lenguaje de Ontologías Web''. Existen muchos lenguajes y herramientas para desarrollar y trabajar con Ontologías, pero ninguna de las existentes hasta el momento resultaba compatible con una arquitectura Web y mucho menos con la Web Semántica. 

El Lenguaje de Ontologías Web, rectifica esto aprovechando una conexión proporcionada por RDF para dar a las Ontologías las siguientes capacidades (W3C, 2005):

\begin{itemize}
\item Capacidad de ser distribuidas y compartidas a través de varios sistemas.
\item Escalable a las necesidades de la Web.
\item Compatible con los estándares internacionales de accesibilidad Web e Internacionalización (I18N).
\item Abierto y extensible.
\end{itemize}

Adicionalmente, OWL es una extensión de RDF Schema para permitir la expresión de relaciones más complejas entre elementos y clases. Algunos de los recursos que tiene OWL y de los que, según el W3C (2005), carece RDFS son las siguientes:

\begin{itemize}
\item Recursos para inferir cuáles elementos que tienen varias propiedades son miembros de una clase en particular.
\item Recursos para determinar si la totalidad de elementos de una clase tendrán una propiedad determinada o si puede ser que sólo algunos elementos la tengan.
\item Recursos para diferenciar relaciones 1:1, 1:N y N:1, permitiendo que las ``claves foráneas'' de las bases de datos puedan ser representadas en la Ontología.
\item Recursos para expresar relaciones entre clases definidas en documentos diferentes a través de la Web.
\item Recursos para definir nuevas clases a partir de uniones, intersecciones y complementos de otras ya existentes.
\item Recursos para restringir rangos y dominios para especificar combinaciones de clases y propiedades.
\end{itemize}

Siendo OWL un lenguaje tan extenso, resulta costoso explotarlo a su máxima expresión en aplicaciones web, ya que resulta muy difícil para el motor de inferencia tomar una decisión rápida respecto a la consulta que se le está realizando, es por ello que OWL se divide en tres (3) subconjuntos, según la extensión y el nivel de detalle y profundidad que se desee dar a la ontología, sin llegar a un modelo de datos en el que la capacidad de decición del razonador no se complique demasiado ni, mucho menos, llegue a ser indecidible. Los subconjuntos de OWL, según Antoniou y Van Harmelen (2003), son los siguientes:

\subsection{OWL Lite}
Es el subconjunto más restrictivo de OWL, pero, aún así, es el más utilizado para la Web Semántica pues incluye los constructores básicos y necesarios para escribir modelos cuya decibilidad no sea muy compleja.
\subsection{OWL DL}
Es un subconjunto de OWL un poco más amplio y muy popular en el ámbito de la Lógica Descriptiva (Descriptive Logic), sigue garantizando eficiencia computacional en la inferencia sobre el modelo restringiendo el uso de algunos constructores de OWL.

\subsection{OWL Full}
Es el conjunto completo del lenguaje OWL, sin ningún tipo de restricción de constructores o utilización de los mismos, la ventaja de OWL Full es que es lo más flexible que existe para el modelado de ontologías para Web Semántica, la desventaja es que el nivel de expresividad del lenguaje es tan rico y poderoso que afecta la decibilidad del razonador.

OWL, se encuentra aún en proceso de desarrollo y de definición como estándar, pero el trabajo se encuentra bien adelantado y es posible desarrollar aplicaciones con lo que existe actualmente. Existen numerosas herramientas que permiten modelar y trabajar con RDF, RDFS y OWL, siendo el más utilizado el editor de ontologías Protegé (desarrollado por la universidad de Stanford), permite editar ontologías en OWL de manera gráfica y, además, tiene varios razonadores (motores de inferencia) integrados para validar la integridad de la ontoligía que se está escribiendo.

\section{SPARQL}

En la Web Semántica, toda la información acerca de los recursos se encuentra modelada en Ontologías definidas mediante RDF, RDFS y OWL (en cualquiera de sus tres dialectos). SPARQL, ``es un lenguaje de consultas para RDF'' (W3C, S/F), es decir, así como en el modelo relacional, pueden consultarse las bases de datos a través de SQL y, de esta manera, obtener subconjuntos de la información almacenada según un conjunto de restricciones, mediante consultas SPARQL puede obtenerse subconjuntos de la información contenida en uno o varios documentos RDF en forma de grafos, dependiendo de las relaciones y de lo que el usuario desee buscar.

SPARQL, viene a ser la evolución de RQL y aún se encuentra en proceso de desarrollo, aunque existen numerosos motores de consulta en SPARQL y se puede trabajar sobre lo que ya existe.

\section{Motores de Inferencia}

Una vez que se tiene la Ontología ya definida, es necesario definir ciertas reglas de inferencia a ser aplicadas sobre ese conocimiento ya descrito. Estas reglas de inferencia son aplicadas por un Motor de Inferencia, que es un programa que explora la base de conocimiento, aplicando ciertas reglas, hasta dar con la solución que mejor se adapte a las necesidades del usuario (Díaz, S/F).

En la Web Semántica, es el motor de inferencia quien traduce la consulta del usuario y evalúa, mediante reglas bien definidas, qué es lo que el usuario realmente está buscando, es por ello que el motor de inferencia juega un papel de suma importancia en este marco, vale acotar, que los motores de inferencia para Web Semántica, deben ser capaces de ``razonar'' sobre RDF, RDFS y OWL.

% TODO: Falta Gestión del conocimiento y procesamiento de lenguaje natural, Servicios Web
\newpage


\chapter{Marco Metodológico}
\label{chap:marcometodologico}

Las metodologías de desarrollo tradicionales, parecieran haber sido diseñadas para proyectos extensos, con equipos numerosos, en los que se goza de roles bien definidos y que, cada rol, cumple con una labor específica en cada etapa del ciclo de vida.

La realidad de este proyecto es otra, sólo se cuenta con una persona para realizar el trabajo de todos los roles, además, el tiempo de entrega es limitado y resulta, además, conveniente entregarlo lo antes posible. Es por ello que se propone trabajar con una metodología Iterativa Incremental bajo el esquema ágil. Los preceptos del esquema de trabajo enmarcado en el Desarrollo Ágil, fueron definidos por Kent Beck año 2001 en un documento denominado El Manifiesto Ágil (The Agile Manifesto), este manifiesto se cita a continuación:

``Estamos descubriendo nuevas maneras de desarrollar software tanto por nuestra propia experiencia como ayudando a terceros. A través de esta experiencia hemos aprendido a valorar:

\begin{itemize}
\item \textbf{Individuos e interacciones} sobre procesos y herramientas.
\item \textbf{Software que funciona} sobre documentación exhaustiva.
\item \textbf{Colaboración con el cliente} sobre negociación de contratos.
\item \textbf{Responder ante el cambio} sobre el seguimiento de un plan.
\end{itemize}

Esto es, aunque los elementos a la derecha tienen valor, nosotros valoramos por encima de ellos a los que están a la izquierda''.

El Trabajo Especial de Grado, tiene un carácter investigativo y aplicativo, es decir, debe investigarse acerca del tema y realizar un pequeño aporte al área de investigación mediante el desarrollo de algún producto final. Por ello, dado el carácter de un TEG, todos y cada uno de los principios definidos en el Manifiesto Ágil de Kent Beck tienen sentido en un proyecto de este tipo pues, debe valorarse ``Individuos e interacciones sobre procesos y herramientas", esto es, debe valorarse más la interacción entre el tesista y el tutor que los procesos y herramientas necesarios para ello, por ejemplo. Debe darse más valor al ``Software de funciona sobre la documentación exhaustiva", en este caso, para la presentación final, debe mostrarse algo funcional, si bien la documentación es importante, no se hará de manera exhaustiva, únicamente lo necesario para poner orden en el proyecto y para que, posteriormente, si el presente trabajo resulta de interés para alguien, pueda entender qué fue lo que se hizo. ``Colaboración con el cliente sobre negociación de contratos", si bien esto no tiene mucha pertinencia al hablar de un TEG, puede interpretarse como que la Escuela de Ingeniería Informática y el Tesista deberían colaborar pues ambos persiguen un objetivo común: innovar. Finalmente, "Responder ante el cambio sobre el seguimiento de un plan" pues al tratarse de un proyecto de investigación, no todos los requerimientos están claramente definidos ni pueden predecirse en su totalidad, es por ello que conforme se va investigando y despejando la incertidumbre en algunos temas, pueden ir surgiendo nuevas cosas que no se esperaban y se debe estar preparado para responder y adaptar el plan a las nuevas condiciones.

Existen numerosas metodologías y ciclos de vida basados en el modelo ágil, para este proyecto en particular se propone utilizar una metodología basada en Scrum para organizar el proyecto y XP (Programación Extrema o eXtreme Programming en inglés) para la organización de cada una de las iteraciones.

\section{Descripción de la Metodología Planteada}

Para describir la metodología que se utilizará para el desarrollo del proyecto, es necesario primero estudiar las dos metodologías antes mencionadas:

\subsection{SCRUM}

SCRUM, segun Roger Pressman (2010), es una metodología creada por Jeff Sutherland y su equipo de desarrollo a principios de la década de 1990. Los principios de SCRUM están enmarcados dentro del Manifiesto Ágil y es un proceso que lleva el Desarrollo de Software a través de las siguientes actividades: requerimientos, análisis, diseño, evolución y entrega. Cada una de esas actividades son realizadas dentro de un patrón de trabajo llamado Sprint, todo el trabajo realizado dentro de un Sprint es adaptado al problema y frecuentemente modificado por el equipo de desarrollo a medida que las condiciones van cambiando.

SCRUM hace énfasis en la utilización de procesos de software que son efectivos en proyectos con tiempos cortos de entrega y requerimientos cambiantes. Esos procesos, en general, definen dos grandes actividades de desarrollo:

\begin{itemize}
\item \textbf{Backlog:} el backlog, constituye una lista priorizada de requerimientos que agregan valor de negocio al producto, estos requerimientos pueden ser agregados en cualquier momento y, de esta manera, se introducen los cambios en el proyecto (Pressman, 2010). El backlog puede ser, además, modificado en cualquier momento para adaptar las prioridades a los cambios del negocio.
\item \textbf{Sprints:} los Sprints constituyen paquetes de trabajo que son necesarios para desarrollar un requerimiento o un conjunto de ellos (Pressman, 2010). Los cambios no pueden ser introducidos dentro de un Sprint, de esta manera se asegura que el trabajo que se realiza es estable, pues los requerimientos que se seleccionan para ser desarrollados en un Sprint ya deben estar definidos y se debe tener la certeza de que, en caso de cambiar, será manejable.
\item \textbf{Scrum Meetings:} son reuniones cortas que, usualmente, se realizan todos los días durante un sprint (Pressman, 2010). Durante los Scrum Meetings se responden a tres preguntas básicas:
\begin{itemize}
\item ¿Qué has hecho desde la última reunión?.
\item ¿Cuáles obstáculos has encontrado?.
\item ¿Cuál es tu planificación hasta la próxima reunión?.
\end{itemize}
\item \textbf{Demos:} se van entregando los resultados de las funcionalidades desarrolladas al cliente para que puedan ser evaluados (Pressman 2010). De esta manera, se va entregando Software Listo y Funcional que agrega valor al negocio del cliente en cada iteración.
\end{itemize}

De esta manera, SCRUM permite que los equipos trabajen de manera eficiente y estable en proyectos en los que la incertidumbre siempre está presente.

\subsection{eXtreme Programming}

Para Sommenvile (2007), eXtreme Programming (XP) es la más utilizada y la más conocida de las metodologías ágiles. Fue inicialmente concebida por Kent Beck a finales de la década de 1980 y se enfoca en hacer énfasis en la etapa de implementación del ciclo de vida, tal como su nombre lo indica, en la parte de programación o codificación, tomando como buena práctica la programación en parejas, de esta manera, mientras uno de los desarrolladores codifica, el otro está observando y advirtiendo de errores a quien escribe el código.

Kent Beck definió cinco valores para establecer las bases de XP como metodología. Cada uno de esos valores se utilizan para poner en marcha las actividades, acciones y tareas de XP (Pressman, 2010):

\begin{itemize}
\item \textbf{Comunicación}: entre los clientes y los desarrolladores. XP valora más la comunicación informal por ser más rápida en comparación con volúmenes enormes de documentación como medio de comunicación principal.
\item \textbf{Simplicidad}: XP restringe a los desarrolladores a realizar el diseño y la implementación de código que satisfaga los requerimientos actuales en vez de pensar en factores futuros, si el diseño puede ser mejorado, puede ser modificado posteriormente.
\item \textbf{Retroalimentación}: durante un proyecto en necesario muchas veces regresar a una etapa previa y modificar algo para que los requerimientos actuales funcionen sin alterar los que ya fueron implementados. El cambio es algo constante en todo proyecto, pero para poder hacerlo correctamente, se necesita retroalimentación.
\item \textbf{Coraje}: implementar código sin miedo a las consecuencias, pero siempre de manera coherente.
\end{itemize}

En eXtreme Programming, no se libera una pieza de software hasta que está totalmente funcional y probada, para asegurar que el software funciona, se desarrollan Pruebas Unitarias. 

La metodología propuesta para el desarrollo del presente proyecto, toma la organización de SCRUM y los valores y la filosofía basada en la codificación de XP para llevar a cabo las tareas y actividades concernientes al desarrollo del producto final de este trabajo, obviando la exigencia de XP acerca de la programación en parejas ya que el presente trabajo es realizado por sólo una persona.

%TODO: Terminar según se vea necesario
\subsection{Plan General de Trabajo}
Se plantea, organizar el desarrollo en las siguientes iteraciones, con una duración de tres a cuatro semanas cada una:

\begin{enumerate}
\item Investigación y análisis.
\item Diseño de la arquitectura del sistema.
\item Establecimiento del ambiente de desarrollo.
\item Anlálisis y desarrollo de la ontología.
\item Integración del modelo de conocimiento con el API seleccionado.
\item Desarrollo del traductor de lenguaje natural a SPARQL.
\item Desarrollo del motor de búsqueda.
\item Desarrollo e integración con las vistas.
\end{enumerate}

\newpage


\chapter{Desarrollo}
\label{chap:desarrollo}

El desarrollo de este trabajo, se describe siguiendo la metodología planteada anteriormente. Las iteraciones citadas en el \textit{Marco Metodológico} serán descritas a continuación de manera consecutiva:

\section{[Iteración 1]: Investigación y análisis}
Durante esta etapa no se desarrolló ningún tipo de software, es por ello que no se verán historias de usuario, diseño, desarrollo ni pruebas unitarias.

\subsection{Desarrollo de las tareas}
Las tareas de esta iteración se basaron principalmente en la investigación y profundización de los temas relacionados a este trabajo, así como la búsqueda de herramientas para el ensamblaje de la plataforma de desarrollo.

\subsubsection{Investigación de los temas relacionados}
La mayor parte de la investigación realizada fue descrita en el \textit{Marco Teórico} del presente trabajo. El tópico central del presente trabajo resulta ser \textit{La Web Semántica}, por ello, es necesario tomar en cuenta todos los conceptos que se desprenden de dicho tópico. Los conceptos más importantes que se desprenden de \textit{La Web Semántica} son citados a continuación:

\begin{itemize}
\item Ontologías y modelos de conocimiento.
\item Motores de inferencia.
\item Servicios Web.
\end{itemize}

Las \textit{Ontologías} y \textit{Modelos de conocimiento}, son el marco central de la \textit{Web Semántica}. Es lo que establece toda la estructura de metadatos en la que se basa el etiquetado e indexado de los recursos que forman parte de la base de conocimientos. Además, establece las relaciones entre las metaentidades que conforman la \textit{Ontología}. Define el vocabulario que modela el dominio del problema a ser resuelto.

Los \textit{Motores de inferencia} son herramientas que, dadas ciertas reglas sobre un \textit{modelo de conocimiento}, son capaces de deducir nueva información y nuevas relaciones entre las metaentidades y, por lo tanto, nuevas relaciones entre los recursos dentro de la base de conocimientos.

Los \textit{Servicios Web}, si bien son un concepto que ya tiene tiempo, cobran especial importancia en la \textit{Web Semántica} pues, cuando se haga efectiva la evolución de la Web 2.0 a la Web 3.0, debe garantizarse total interoperabilidad e independencia de plataformas para que los agentes puedan ser capaces de consultar e intercambiar información y \textit{modelos de conocimiento} de distintas fuentes. Los \textit{servicios} web exponen toda la funcionalidad de un sistema a través de métodos invocables utilizando HTTP sobre TCP. Es por ello que, a través de \textit{Servicios Web}, es posible desarrollar servicios sobre cualquier plataforma y, sin importar en cual otra esté desarrollado, cualquier cliente será capaz de consumir los servicios pues la invocación se encapsula dentro de un protocolo común entre ambas plataformas.

\subsubsection{El problema y los requerimientos}
Del problema general del proyecto, puede plantearse de la siguiente manera: \textit{¿Cómo facilitar el acceso a la información acerca de Ciencias de la Computación a personas interesadas en el área}. De este problema, puede identificarse los siguientes requerimientos:

\begin{itemize}
\item Una interfaz web donde los usuarios puedan interactuar con el sistema.
\item Una interfaz de edición que permita a usuarios autorizados agregar nuevos recursos y extender la base de conocimiento.
\item Una interfaz de edición que permita a usuarios autorizados agregar metainformación nueva y extender el modelo de conocimiento.
\item Un componente de traducción que convierta consultas en lenguaje natural a SPARQL, el lenguaje de consultas sobre RDF/RDFS/OWL.
\item Desarrollo de un modelo de conocimiento del área de Ciencias de la Computación e Ingeniería Informática.
\item Un componente capaz de realizar consultas e inferencias sobre el modelo de conocimiento realizado.
\end{itemize}

\subsubsection{Selección de la plataforma}
Al ser una aplicación bajo la arquitectura Cliente-Servidor, hay que tener especial cuidado en seleccionar las mejores herramientas para desarrollar cada uno de los componente de dicha arquitectura.

Al principio, se planteó llevar a cabo todo el desarrollo utilizando \textit{Python} como lenguaje de programación y las librerías nativas disponibles a través de \textit{easy\_install}, pues \textit{Python} es uno de los pocos lenguajes que ofrece librerías nativas para la manipulación de documentos RDF, además de las ventajas propias que ofrece el lenguaje desde el punto de vista sintáctico y de facilidad de aprendizaje y de programación. Pero a la hora de buscar Frameworks que permitieran el desarrollo de aplicaciones basadas en Web Semántica y Motores de Inferencia, la información y la cantidad de herramientas resultó limitada.

Debido a la situación expuesta en el párrafo anterior, se planteó evaluar el desarrollo por separado: seleccionar la mejor plataforma para desarrollar el servidor y, por otra parte, la mejor plataforma para desarrollar el cliente:

\subsubsection{Servidor}


\subsubsection{Selección de los frameworks de desarrollo}


\subsubsection{Selección del medio de persistencia}


\newpage

\chapter{Resultados}
\label{chap:resultados}

Luego de haber finalizado el desarrollo y la validación del prototipo, se presentan a continuación los resultados obtenidos:

\begin{itemize}
    \item Se obtuvo una ontología escrita en \textit{OWL} aplicable para modelar cualquier Disciplina de las ciencias o las humanidades a través de la creación de las instancias correspondientes a ella.
    \item Se obtuvo un modelo de conocimiento del área de \textit{Ciencias de la Computación} escrito en \textit{OWL}, basado en la ontología mencionada en el punto anterior.
    \item Se obtuvo una herramienta de carácter prototípico para la administración de las instancias (anotaciones) del modelo de conocimiento mencionado en el punto anterior.
    \item Se obtuvo un lenguaje básico de consultas basado en \textit{Forma Normal Disyuntiva}.
    \item Se obtuvo un buscador sobre el modelo de conocimiento del área de \textit{Ciencias de la Computación} que implementa el lenguaje mencionado en el punto anterior.
\end{itemize}

\chapter{Conclusiones y Recomendaciones}
\label{chap:conclusiones}

\section{Conclusiones}
...

\section{Recomendaciones}

...

\pagestyle{plain}

\begin{thebibliography}{20}

    \bibitem{swp} Antoniou, G. y Van Harmelen, F. (2003). \emph{A Semantic Web Primer}. Massachusetts: MIT Press.

    \bibitem{planeta} Cobo, C. y Pardo, H. (2007). \emph{Planeta 2.0: Inteligencia Colectiva o Medios Fast Food}. México DF: Grup de Recerca d'Interaccions Digitals.

    \bibitem{exito} Éxito Explorador (Agosto 31, 2010). \emph{Estadísticas Mundiales de Internet} [Datos en línea]. Disponible: http://www.exitoexportador.com/stats.htm [Consulta: 2010, Noviembre 28].

    \bibitem{gbws} W3C. \emph{Guía Breve de la Web Semántica} [Documento en línea]. Disponible: http://www.w3c.es/divulgacion/guiasbreves/websemantica [Consulta: 2010, Noviembre 28].

    \bibitem{ricardo} Casanova, R. (2010). \emph{El Modelo Web 2.0}. Presentación sobre el modelo de negocios orientado a la Web 2.0.

    \bibitem{rdfp} W3C. \emph{RDF Primmer} [Documento en línea]. Disponible: http://www.w3.org/TR/2004/REC-rdf-primer-20040210/ [Consulta: 2010, Noviembre 28].

    \bibitem{rdf} W3C. \emph{Resource Description Framework (RDF)} [Documento en línea]. Disponible: http://www.w3.org/RDF/ [Consulta: 2010, Noviembre 29].

    \bibitem{rdffaq} GSI. \emph{RDF Schema (RDFS)} [Presentación en línea]. Disponible: www.gsi.dit.upm.es/~gfer/ssii/RDFS.pdf [Consulta: 2010, Noviembre 30].

    \bibitem{owlfaq} W3C. \emph{Preguntas Frecuentes del OWL} [Documento en línea]. Disponible: http://www.w3c.es/Traducciones/es/SW/2005/owlfaq [Consulta: 2010, Noviembre 30].

    \bibitem{owl} W3C. \emph{OWL: Use Cases and Requirements} [Documento en línea]. Disponible: http://www.w3.org/TR/2004/REC-webont-req-20040210/\#onto-def [Consulta: 2010, Diciembre 01].

    \bibitem{sparql} W3C. \emph{SPARQL Query Language for RDF} [Documento en línea]. Disponible: http://www.w3.org/TR/rdf-sparql-query/ [Consulta: 2010, Diciembre 02].

    \bibitem{se} Diaz, S. (s/f). \emph{Sistemas Expertos. Un paso en la simulación del razonamiento humano} [Documento en línea]. Disponible: http://www.monografias.com/trabajos23/sistemas-expertos/sistemas-expertos.shtml [Consulta: 2010, Noviembre 30]

    \bibitem{manifiesto} Beck, K. et al. (2001). \emph{Manifesto for Agile Software Development} [Documento en línea]. Disponible: http://agilemanifesto.org/ [Consulta: 2010, Diciembre 02].

    \bibitem{pressman} Pressman, R. (2010). \emph{Software Engineering: a Practitioner's Approach}. New York: McGraw Hill.

    \bibitem{sommervile} Sommervile, I. (2007). \emph{Software Engineering}. New York: Pearson Education.

    \bibitem{foaf} The Friend of a Friend Project. \emph{FOAF Project} [Documento en línea]. Disponible: http://www.foaf-project.org [Consulta: 2011, Enero 15].

    \bibitem{python} Python Software Foundation. \emph{About Python} [Documento en línea]. Disponible en: http://www.python.org/about [Consulta: 2011, Agosto 08].

    \bibitem{java-joyanes} Joyanes, L. y Zahonero I. (2002). \emph{Programación en Java 2}. Madrid: McGraw Hill.

    \bibitem{jena} Jena Community (s/f). \emph{About Jena} [Documento en línea]. Disponible: http://jena.sourceforge.net/index.html [Consulta: 2011, Agosto 08].

    \bibitem{sesame} Autor desconocido (s/f). \emph{About Sesame} [Documento en línea]. Disponible: http://www.openrdf.org/about.jsp [Consulta: 2011, Agosto 08].

    \bibitem{redland} RedLand Community (s/f). \emph{RedLand FAQ} [Documento en línea]. Disponible: http://librdf.org/FAQS.html [Consulta: 2011, Agosto 08].

    \bibitem{cubicweb} CubicWeb Community (s/f). \emph{The Semantic Web is a Construction Game} [Documento en línea]. Disponible http://www.cubicweb.org/ [Consulta: 2011, Agosto 08].

    \bibitem{virtuoso} Virtuoso Community (s/f). \emph{Virtuoso FAQ} [Documento en línea]. Disponoble: http://virtuoso.openlinksw.com/dataspace/dav/wiki/Main/VOSVirtuoso6FAQ [Consulta: 2011, Agosto 08].

    \bibitem{debian} Debian Community (s/f). \emph{About Debian} [Documento en línea]. Disponible: http://debian.org [Consulta: 2011, Diciembre 20]

    \bibitem{techradar} Autor Desconocido (s/f). \emph{10 Best Linux Distros for 2011} [Documento en línea]. Disponible: http://www.techradar.com/news/software/operating-systems/10-best-linux-distros-for-2011-704584?artc\_pg=2 [Consulta: 2012, Enero 10]

    \bibitem{norvig-russel} Norvig, P. y Russell, S. (2004). \emph{Inteligencia Artificial Un Enfoque Moderno} (2a. ed). Madrid: Prentice Hall.

    \bibitem{aho} Aho, A. et al. (2007). \emph{Compilers Principles, Techniques & Tools} (2a. ed). Boston: Pearson Education.

    \bibitem{lucia} Cardoso, L. (2006). \emph{Sistemas de Base de Datos II, Teoría aplicada para profesores y estudiantes}. Caracas: Publicaciones UCAB.

    \bibitem{papelLuque} Luque, M. et al. (s/f). \emph{Caracterización Automática de Perfiles de Usuarios basados en Consultas Lingüísticas Multigranulares usando un Algoritmo Genético Multiobjetivo}. España.

\bibitem{ontoclean} Guarino, N. y Welty, C. (s/f). \emph{An Overview of OntoClean}. Italia.

\bibitem{kuster} Kúster, W y Rauseo, O. (2009). \emph{Implementar un servicio web que permita crear mashups públicos o privados para ser usados a través de la mensajería de texto celular}. Trabajo Especial de Grado para optar al título de Ingeniero en Informática, Universidad Católica Andrés Bello. Caracas.

\bibitem{ieeeacm} IEEE-CS y ACM (2001). \emph{Computing Curricula 2001 Computer Science} [Documento en línea]. Disponible: http://www.acm.org/education/education/education/curric\_vols/cc2001.pdf [Consulta: 2011, Marzo 15].

\bibitem{acm} IEEE-CS y ACM (2004). \emph{Computer Engineering 2004: Curriculum Guidelines for Undergraduate Degree Programs in Computer Engineering} [Documento en línea]. Disponible: http://www.acm.org/education/education/curric\_vols/CE-Final-Report.pdf [Consulta: 2011, Marzo 15].

\bibitem{ieee} IEEE-CS y ACM (2008). \emph{Computer Science 2008: an Interim Revision of CS 2001} [Documento en línea]. Disponible: http://www.acm.org//education/curricula/ComputerScience2008.pdf [Consulta: 2011 Marzo, 15]

\bibitem{progsemweb} Hebeler, J. et al. (2009). \emph{Semantic Web Programming}. Indiana: Wiley Publishing, Inc.

\end{thebibliography}



% es un cap'itulo sin n'umero.
\chapter*{Lista de abreviaturas}  
\label{chap:abrev}
\addcontentsline{toc}{chapter}{Lista de abreviaturas} 


\begin{center}
\begin{longtable}{r p{13.5cm}} % Se usa longtable porque posiblemente no quepa en una sola p'agina

\textbf{3G} & Del espa\~nol: tercera generaci\'on de transmisi\'on de voz y datos a trav\'es de telefon\'ia m\'ovil\\

\textbf{AVD} & Del ingl'es:\textit{ Android Virtual Devices} , en espa\~nol: Dispositivos Virtuales de Android\\

\textbf{ADT} & Del ingl'es:\textit{ Android Development Tools} , en espa\~nol: Herramientas de Desarrollo de Android\\

\textbf{API} & Del ingl'es: \textit{Application Programming Interface}, en espa\~nol: Interfaz de Programaci'on de Aplicaciones.\\

\textbf{BTS} & Del ingl'es: \textit{Base Transceiver Station}, en espa\~nol: Estaci\'on Base.\\

\textbf{BSC} & Del ingl'es: \textit{Base Station Controller}, en espa\~nol: Controlador de Estaciones Base.\\

\textbf{CCS} & Del ingl'es:\textit{ Cascading Style Sheets}, en espa\~nol: Hojas de Estilo en Cascada.\\

\textbf{CMS} & Del ingl'es: \textit{Content Management System}, en espa\~nol: Sistema Manejador de Contenido.\\

\textbf{EDGE} & Del ingl'es: \textit{Enhanced Data rates for GSM of Evolution}, en espa\~nol: Tasas de Datos Mejoradas para la Evoluci\'on de GSM.\\

\textbf{SDK} & Del Ingl'es: \textit{Software Development Kit}, en espa\~nol: Kit de Desarrollo de Software.\\

\textbf{GNU} & Del ingl'es:  GNU \textit{is Not} Unix, en espa\~nol: Acr'onimo recursivo que significa GNU No es Unix. \\

\textbf{GPL} & Del ingl'es: \textit{General Public License}, en espa\~nol: Licencia P'ublica General. \\

\textbf{GPS} & Del ingl'es: \textit{Global Positioning System}, en espa\~nol: Sistema de Posicionamiento Global. \\

\textbf{GPRS} & Del ingl'es: \textit{General Packet Radio Service}, en espa\~nol: Servicio General de Paquetes V\'ia Radio. \\

\textbf{HTTP} & Del ingl'es: \textit{HyperText Transfer Protocol}, en espa\~nol: Protocolo de Transferencia de Hipertexto. \\

\textbf{IDE} & Del ingl'es:  \textit{integrated development environment}, en espa\~nol: Entorno de desarrollo integrado. \\

\textbf{JDK}  & Del ingl'es: \textit{Java Development Kit}, en espa\~nol:Kit de desarrollo de Java.\\

\textbf{jSON}  & Del ingl'es: \textit{JavaScript Object Notation}, en espa\~nol:Kit Notaci\'on de Objetos de JavaScript.\\

\textbf{OMS} & Del Espa\~nol:  Organizaci\'on Mundial de la Salud.\\

\textbf{RAE} & Del Espa\~nol: Real Academia Espa\~nola. \\

\textbf{RF} & Del espa\~nol: Ruta de Radio-Fecuencia\\

\textbf{TIC} & Del espa\~nol: Tecnolog\'ias de la informaci\'on y la comunicaci\'on\\

\textbf{OS} & Del ingl'es: \textit{Operating System}, en espa\~nol: Sistema Operativo. \\

\textbf{PHP} & Del ingl'es: \textit{PHP Hypertext Pre-processor}.\\

\textbf{POI} & Del ingl'es: \textit{Point of Interest}, en espa\~nol: \textit{Puntos de Inter\'es}.\\

\textbf{XML} & Del ingl'es: \textit{Extensible Markup Language}, en espa\~nol: Lenguaje de Marcado Extensible. \\

\textbf{XP} & Del ingl'es: \textit{Xtreme Programming}, en espa\~nol: Programaci\'on extrema. \\

\end{longtable}\end{center}

\addtocontents{toc}{\vspace{2em}} % Add a gap in the Contents, for aesthetics

\pagenumbering{roman}

\pagestyle{empty}

\vspace*{3.5cm}
\begin{center}
\begin{LARGE}
Aplicaci\'on de apoyo a los procesos de ubicaci\'on de nuevos enlaces corporativos y de selecci\'on de los sitios de transmisi\'on a los que se conectar\'an

(AP\'ENDICES)
\end{LARGE}
\vspace{\fill}
{
\begin{tabular}{p{6.5cm}l}
& \textbf{Ricardo Recaredo}\tabularnewline
 & \textbf{Viviana Trujillo}\tabularnewline
\multicolumn{1}{l}{} & \multicolumn{1}{l}{}\tabularnewline
\multicolumn{1}{l}{}& \multicolumn{1}{l}{\textbf{Jorge Robaina}}\tabularnewline
\multicolumn{1}{l}{} & \multicolumn{1}{l}{}\tabularnewline
\multicolumn{1}{l}{} & \multicolumn{1}{l}{}\tabularnewline
\multicolumn{1}{l}{}& \multicolumn{1}{l}{Caracas, 29 de septiembre de 2010} \tabularnewline
\end{tabular}
}
\end{center}


\pagenumbering{arabic}
\setcounter{page}{1} 

\pagestyle{ruled}

\appendix % Cue to tell LaTeX that the following 'chapters' are Appendices
\clearpage

%Marco Teorico
\include{apendiceG} 		% Appendix GSM
\chapter{M\'etodo de Ponderaci\'on Lineal (scoring)}
\label{AppendixMPL}

\section{Matriz de ponderaci\'on del Score para el m\'etodo de ponderaci\'on lineal}
Para tener una visi\'on mas amplia del escenario de evaluaci\'on y de la comparativa de los Scores sobre los diferentes criterios para cada una de las alternativas, es recomendable elaborar una Matriz de ponderaci\'on de alternativas basada en el peso asignado a los criterios. 

Lo usual es que en la Matriz De Ponderaci\'on en su forma general[\ref{tbl:matrizPonderacion}], en la primera columna se presenten las alternativas a ser evaluadas y en las siguientes columnas los criterios, dejando la primera fila para identificar los respectivos criterios y los rangos de sus pesos y las restantes casillas de la matriz para realizar la valoraci\'on propiamente dicha, y se conserva la \'ultima columna para completar la evaluaci\'on de cada alternativa, sumando los puntos acumulados por la misma, en su respectiva fila.

\renewcommand{\baselinestretch}{1}

\nuevatabla{matrizPonderacion}
{\footnotesize Matriz De Ponderaci\'on en su forma m\'as general}
{\tiny Fuente: Matrices de Ponderaci\'on para la Evaluaci\'on de Proveedores. 2007}
{p{2cm}p{1.5cm}p{1.5cm}p{0.5cm}p{2cm}p{1.5cm}p{1cm}}
{
\multicolumn{1}{>{\columncolor{vivinicered}}c}{} &
\multicolumn{1}{>{\columncolor{vivinicered}}c}{\footnotesize \textcolor{white}{\textbf{Criterio 1}}} &
\multicolumn{1}{>{\columncolor{vivinicered}}c}{\footnotesize \textcolor{white}{\textbf{Criterio 2}}}& 
\multicolumn{1}{>{\columncolor{vivinicered}}c}{}  & 
\multicolumn{1}{>{\columncolor{vivinicered}}c}{\footnotesize \textcolor{white}{\textbf{Criterio n-1}}} & 
\multicolumn{1}{>{\columncolor{vivinicered}}c}{\footnotesize \textcolor{white}{\textbf{Criterio n }}} & 
\multicolumn{1}{>{\columncolor{vivinicered}}c}{}\\
 
\multicolumn{1}{>{\columncolor{vivinicered}}c}{\footnotesize \textcolor{white}{\textbf{Alternativas}}} & 
\multicolumn{1}{>{\columncolor{vivinicered}}c}{\footnotesize \textcolor{white}{\textbf{Peso x a X}}} & 
\multicolumn{1}{>{\columncolor{vivinicered}}c}{\footnotesize \textcolor{white}{\textbf{Peso y a Y}}} & 
\multicolumn{1}{>{\columncolor{vivinicered}}c}{\footnotesize \textcolor{white}{\textbf{\ldots{}}}} & 
\multicolumn{1}{>{\columncolor{vivinicered}}c}{\footnotesize \textcolor{white}{\textbf{Peso w a W}}}  & 
\multicolumn{1}{>{\columncolor{vivinicered}}c}{\footnotesize \textcolor{white}{\textbf{Peso z a Z}}} & 
\multicolumn{1}{>{\columncolor{vivinicered}}c}{\footnotesize \textcolor{white}{\textbf{Total}}} \\ \hline

\multicolumn{1}{c}{\footnotesize Alternativa 1} &
\multicolumn{1}{c}{\footnotesize P1,1} & 
\multicolumn{1}{c}{\footnotesize P1,2} & 
\multicolumn{1}{c}{\footnotesize \ldots{}} & 
\multicolumn{1}{c}{\footnotesize P1,n-1}  & 
\multicolumn{1}{c}{\footnotesize P1,n} & 
\multicolumn{1}{c}{\footnotesize Total 1} \\ 

\multicolumn{1}{>{\columncolor{cap}}c}{\footnotesize Alternativa 2} & 
\multicolumn{1}{>{\columncolor{cap}}c}{\footnotesize P2,1} & 
\multicolumn{1}{>{\columncolor{cap}}c}{\footnotesize P2,2} & 
\multicolumn{1}{>{\columncolor{cap}}c}{\footnotesize \ldots{}} & 
\multicolumn{1}{>{\columncolor{cap}}c}{\footnotesize p2,n-1}  & 
\multicolumn{1}{>{\columncolor{cap}}c}{\footnotesize P2,n} & 
\multicolumn{1}{>{\columncolor{cap}}c}{\footnotesize Total 2}\\

\multicolumn{1}{c}{\footnotesize \ldots{}} &
\multicolumn{1}{c}{\footnotesize \ldots{}} & 
\multicolumn{1}{c}{\footnotesize \ldots{}} & 
\multicolumn{1}{c}{\footnotesize \ldots{}} & 
\multicolumn{1}{c}{\footnotesize \ldots{}} & 
\multicolumn{1}{c}{\footnotesize \ldots{}} & 
\multicolumn{1}{c}{\footnotesize \ldots{}}\\

\multicolumn{1}{>{\columncolor{cap}}c}{\footnotesize Alternativa m-1} &  
\multicolumn{1}{>{\columncolor{cap}}c}{\footnotesize Pm-1,1} & 
\multicolumn{1}{>{\columncolor{cap}}c}{\footnotesize Pm-1,2} & 
\multicolumn{1}{>{\columncolor{cap}}c}{\footnotesize \ldots{}} & 
\multicolumn{1}{>{\columncolor{cap}}c}{\footnotesize Pm-1,n-1} & 
\multicolumn{1}{>{\columncolor{cap}}c}{\footnotesize Pm-1,n} & 
\multicolumn{1}{>{\columncolor{cap}}c}{\footnotesize Total m-1}\\

\multicolumn{1}{c}{\footnotesize Alternativa m} & 
\multicolumn{1}{c}{\footnotesize Pm,1} & 
\multicolumn{1}{c}{\footnotesize Pm,2} & 
\multicolumn{1}{c}{\footnotesize \ldots{}} & 
\multicolumn{1}{c}{\footnotesize Pm,n-1} & 
\multicolumn{1}{c}{\footnotesize Pm,n} & 
\multicolumn{1}{c}{\footnotesize Total m}\\ \hline
}

\renewcommand{\baselinestretch}{1.5}

En la forma m\'as usual los rangos son todos iguales, es decir x = y = \ldots{} = w = z, en general igual a cero (0) o uno (1) como valor inicial y X = Y = \ldots{} = W = Z = a un valor fijo, generalmente diez (10), veinte (20), cincuenta (50), cien (100) o mil (1000), de acuerdo a lo dif\'icil que puede ser diferenciar entre una alternativa y otra. La otra variante, comnmente usada da valores diferentes a los rangos de cada criterio, en este caso los valores de x, y, w, z, no necesariamente son ceros o unos, y los valores de X, Y, W, Z, suelen ser diferentes, entendi\'endose que entre mayor sea este \'ultimo valor, mayor peso se le desea dar a ese criterio y entre mayor sea la diferencia del mayor valor de un criterio menos el menor valor del mismo (X \textendash{} x, \ldots{}, Z \textendash{} z), se desea ser m\'as diferenciador en cuanto a este criterio. 

 		% Appendix Metodo de ponderacion lineal

\addtocontents{toc}{\vspace{2em}}  

%\backmatter

\end{document}